% !PRINTED 360
\weberpage{360}

This gives the development of the fraction $m:m_1$ or $m:n$ into a \gesperrt{continued fraction}. We obtain successively
\weq{5}{
\frac{m}{n} = a_0 + \cfrac{1}{\dfrac{m_1}{m_2}}, \quad
\frac{m}{n} = a_0 + \cfrac{1}{a_1 + \cfrac{1}{\dfrac{m_2}{m_3}}} \cdots,
}
and finally
\weq{6}{
\frac{m}{n} = a_0 + \cfrac{1}{a_1 + \cfrac{1}{a_2 + \cdots + \cfrac{1}{a_{\nu-1} + \cfrac{1}{m_\nu}}}},
}
for which, for the sake of more convenient printing, we also
\weq{7}{
\begin{aligned}
\frac{m}{n} &= \left(a_0, \frac{m_1}{m_2}\right), \quad \frac{m}{n} = \left(a_0, a_1, \frac{m_2}{m_3}\right)\\
\frac{m}{n} &= (a_0, a_1, a_2 \ldots a_{\nu-1}, m_\nu)
\end{aligned}
}
write.

To this the following remark must still be added: every positive rational fraction can be put into the form $m:n$ in such a way that $m$ and $n$ are relatively prime, and hence we can develop every such fraction into a continued fraction of the form (6) or (7). This also holds for negative rational fractions, if we admit negative values for the number $a_0$ as well. The $a_1, a_2 \ldots a_{\nu-1}, m_\nu$ are, however, positive integers even then. These numbers are completely and unambiguously determined. Only in one respect is an arbitrariness still open to us.

Namely, according to the determination made up to this point, the last of the partial denominators $m_\nu$ is greater than 1, since we have assumed that $m_{\nu+1}$ was to be the first among the residues that is equal to 1. Accordingly, if $a_\nu$ is a positive integer, we can set either
\[
m_\nu = a_\nu,
\]
or
\[
m_\nu = a_\nu + \frac{1}{1}
\]
and therefore obtain the two representations by continued fractions:
%CONT>>
